\section{Discussion}
\label{sec:discussion}

\subsection{Strengths}
\label{sec:strengths}

Our approach offers several advantages for clinical deployment. \textbf{Zero forgetting} eliminates the need for replay buffers or complex regularization; adapter isolation guarantees BWT=0. \textbf{Parameter efficiency} ($<$0.1\% per task) enables fast adaptation and low storage. \textbf{Modularity} allows adding new tasks without retraining prior adapters. The frozen backbone preserves pretrained representations, which is critical when downstream data is scarce. Such properties are desirable when a clinical site incrementally adds analysis capabilities (e.g., brain age estimation) to an existing pipeline with few new annotations and without retaining prior-task data.

\subsection{Limitations}
\label{sec:limitations}

\textbf{T1 performance gap:} LoRA (0.60 Dice) underperforms Sequential Linear (0.79) and Sequential FT (0.80) on segmentation. This may stem from LoRA capacity or decoder adaptation; ablations show encoder+decoder LoRA is necessary. \textbf{T2 systematic bias:} Wilcoxon test on brain age residuals shows significant systematic underestimation ($p<0.001$), possibly due to few-shot training, IXI demographics, and the imputation of missing ages to 50.0\,yr. \textbf{Infarct detection excluded:} The FOMO formulation includes infarct detection (ISLES 2022) as a binary classification task, but we did not evaluate it. ISLES contains mostly infarct-positive cases; a balanced classifier requires healthy (negative) controls, which ISLES does not provide. Supplementing with IXI healthy subjects would enable evaluation; we defer this to future work. \textbf{T2 MAE overfitting (Sequential FT, EWC):} Both Sequential FT (MAE 0.005) and EWC (MAE 0.001) achieve implausibly low T2 MAE; per-task linear probing on full eval achieves 0.063, suggesting these methods overfit on the few-shot validation set. We flag both in Table~\ref{tab:main}. \textbf{IXI age imputation:} Missing ages in IXI metadata are imputed to 50.0\,yr; the six ``best'' T2 samples (lowest error) in the appendix all have GT=50.0\,yr, indicating a concentration of imputed labels that limits validation interpretability. \textbf{Task order:} We evaluated T2$\to$T3 (Phase 1--2) and T3$\to$T2 (Phase 3, Table~\ref{tab:t32}). In T3$\to$T2, Sequential FT shows severe T3 forgetting (BWT$\approx$7.2); LoRA maintains BWT=0. n\_shot=64 results (EWC, LwF) are in Table~\ref{tab:shot64}.

\subsection{Related Work}
\label{sec:related}

Foundation models for brain MRI~\cite{fomo2025,unetr} and PEFT (LoRA~\cite{hu2022lora}, Adapters~\cite{houlsby2019}) are well-established. Continual learning in medical imaging~\cite{kirkpatrick2017,li2017lwf} typically addresses class-incremental segmentation. Our work combines frozen FOMO-style backbones with LoRA for \emph{task-incremental} few-shot continual learning across heterogeneous outputs (segmentation and regression), which is less explored.
